%% Introduction. Sources: references/README.md (Shananin 1997a/b, 1999; Agaltsov–Molchanov–Shananin 2018;
%% Agaltsov 2015, 2016; Molchanov 2013a/b; Kuosmanen–Post–Scholtes 2007; Apesteguia–Ballester 2010).
\section{Introduction}\label{sec:intro}

\paragraph{The model.}
An industry produces a homogeneous output from $d$ inputs. Following \citet{Houthakker1955} and \citet{Johansen1972}, the industry is a continuum of production units. A unit is described by its \emph{technology} $x\in\Rpp[d]$, the vector of input requirements per unit of output, and the units are distributed over technologies by a nonnegative measure $\mu$ on $\Rp[d]$, the \emph{distribution of capacities}: $\mu(X)$ is the total capacity of the units with technologies in $X$. In the original model a unit cannot substitute its inputs, and its cost per unit of output at input prices $p\in\Rpp[d]$ is $\ip px$. \citet{Sato1975} asked for a version of the model in which inputs substitute at the level of a single unit; it was built and studied by \citet{Shananin1997a,Shananin1997b}. There every unit has the same production function $F_0(u^1/x^1,\dots,u^d/x^d)$ of its input flows $u$, and $F_0$ is \emph{neoclassical}: nonnegative, nondecreasing, concave, continuous and positively homogeneous of degree one. The cost of a unit of output by the technology $x$ is then $h(p\circ x)$, where $p\circ x$ is the coordinatewise product and $h$ is the Young transform of $F_0$, $h(q)=\inf\set{\ip qu/F_0(u)\mid u\ge0,\ F_0(u)>0}$ \citep[Lemma~4.1]{Shananin1997a}. The function $h$ is again neoclassical; we call it the \emph{unit cost function} and work with it throughout.

When the output price is $p_0$ and the input prices are $p$, the units with $h(p\circ x)<p_0$ are profitable and run at full capacity, the units with $h(p\circ x)>p_0$ are idle \citep[Thm.~4.1]{Shananin1997a}. Hence the output of the industry is
\begin{equation}\label{eq:intro-output}
y=\mu\set{x\in\Rpp[d]\mid h(p\circ x)<p_0}.
\end{equation}
Only prices and outputs are observed; the distribution $\mu$, the input flows of the units and the function $h$ are latent. Since $h$ is homogeneous, \eqref{eq:intro-output} depends on the prices through $p/p_0$ only, and we normalize the output price to one. The data are $T$ observations $(p_t,y_t)$, $p_t\in\Rpp[d]$, $y_t\ge0$, and the question is whether they are generated by \eqref{eq:intro-output} for some neoclassical $h$ and some $\mu$.

\paragraph{The inverse problem with complete data.}
The macro-level description of the industry is its profit function
\begin{equation}\label{eq:intro-profit}
\Pi(p_0,p)=\int_{\Rp[d]}\bigl(p_0-h(p\circ x)\bigr)_+\,\mu(dx),
\end{equation}
whose derivative in $p_0$ is the output \eqref{eq:intro-output}, and whose second derivative is the generalized Radon transform of $\mu$ over the level hypersurfaces $\set{x\mid h(p\circ x)=p_0}$ \citep{Shananin1997b,AgaltsovMolchanovShananin2018}. Recovering $\mu$ from $\Pi$ is an inverse problem of integral geometry, and its solvability criteria have a common form: $\Pi$ is of the form \eqref{eq:intro-profit} with a nonnegative $\mu$ if and only if an explicit linear transform of the data is completely monotone, that is, by Bernstein's theorem, the Laplace transform of a nonnegative measure. For linear $h$, the original Houthakker--Johansen model, the transform is $G(s)=\int_0^\infty e^{-\tau}\,d_\tau\,\partial_\tau\Pi(s,\tau)$, and the criterion reads $(-1)^kD_{\xi_1}\cdots D_{\xi_k}G(\lambda s)\ge0$ for all $k$ and all directions $\xi_i$ in an open cone, together with homogeneity and a growth condition \citep{HenkinShananin1990}. For an arbitrary positively homogeneous $h$ the data are first divided, in the Mellin domain, by the symbol $\rho_h(z)=\Gamma(z_1)\cdots\Gamma(z_d)\Gamma(2+z_1+\dots+z_d)/\mathcal M[e^{-h}](z)$, and the quotient must be completely monotone \citep{Agaltsov2015}; for a CES cost function the inversion is explicit \citep[Prop.~4.1]{AgaltsovMolchanovShananin2018}. The distribution $\mu$ is unique in $L^r$ if and only if the zero set of $\mathcal M[e^{-h}]$ on the plane $\operatorname{Re}z=c$ is meagre, which holds for CES and nested CES cost functions \citep{Agaltsov2016}; the proof rests on Wiener's Tauberian theorem for the multiplicative shifts of the kernel. In this regime the cost function $h$ enters analytically, through its Mellin symbol, and is assumed known.

\paragraph{The inverse problem with finite data and a fixed cost function.}
Statistical agencies publish prices and outputs for finitely many periods, not the profit function. For $T$ observations and a given $h$, \citet{Shananin1999} posed the identification problem as the moment problem \eqref{eq:intro-output}, $t\in[T]$, and proved that it is solvable if and only if $y$ lies in the cone spanned by the $0/1$ vectors recording, for each cell of the partition of $\Rpp[d]$ by the hypersurfaces $h(p_t\circ x)=1$, in which periods the cell is profitable (\cref{lem:signature} below). The combinatorics of this cone was studied by \citet{AgaltsovMolchanovShananin2018} and \citet{Molchanov2013a,Molchanov2013b}: the cone has integer facet normals; the inequalities $\sum_{t\in\Omega_+}y_t\ge\sum_{t\in\Omega_-}y_t$ with coefficients $0,\pm1$ that hold on all cells are necessary for solvability, and they are sufficient if and only if the cone is discretely convex, which for $d=2$ fails already at $T=6$; for a parametric family such as CES the admissible parameters are found by tracking the partition, which for two inputs has $O(T^2)$ cells and changes at $O(T^3)$ parameter values, so that the test is a linear program of polynomial size. The question of which families of cells are realizable by straight lines leads to the stretchability of pseudoline arrangements, which is NP-hard \citep{Shor1991,Molchanov2013a}.

\paragraph{This paper: finite data and an unknown cost function.}
We treat the third regime: finitely many observations and an $h$ ranging over all neoclassical functions, so that the question is whether the data are compatible with the generalized Houthakker--Johansen model at all, with $h$ and $\mu$ both latent. The data are $T$ samples $y_t=\partial_{p_0}\Pi(1,p_t)$ of the derivative of the profit function, and the parallels with the two regimes above organize the paper.
\begin{itemize}
\item \emph{Cone duality.} In all three regimes the data lie in the cone generated by the elementary responses of single technologies, and the criterion is a description of that cone. With complete data the elementary responses are the kernels $(p_0-h(p\circ x))_+$ and the description is Bernstein's: infinitely many linear inequalities, the alternating signs of all derivatives. With finite data the responses are the sampled indicators $\ind S$, $S$ the set of periods in which a cell is profitable, and the description is a finite system of linear inequalities (\cref{thm:solvability,cor:farkas}).
\item \emph{The cost function drops out.} Finitely many prices see only the combinatorics of the partition, not the analytic kernel. The union over all neoclassical $h$ of the families of cells is the family of \emph{reachable spectra}, defined by linear separation alone (\cref{def:reachable}), and one piecewise linear $h$ realizes all of them (\cref{lem:patch}). Hence the criterion for an unknown $h$ is as explicit as Shananin's criterion for a known one, whereas with complete data the division by $\rho_h$ cannot be avoided and the range of \eqref{eq:intro-profit} over all $h$ is not known.
\item \emph{No finite Bernstein theorem.} The finite analogues of complete monotonicity are the inequalities with coefficients $0,\pm1$: dominance $y_t\le y_s$ for $p_s<p_t$ is the first difference, and the \emph{cover inequalities} $y_t\le\sum_{s\in A}y_s$, for $p_t$ above the convex hull of $\set{p_s}_{s\in A}$, are the higher ones (\cref{lem:cover,prop:cover}). They describe the generators of the cone exactly, but not the cone: for a fixed $h$ this is the failure of discrete convexity found by Molchanov, and for an unknown $h$ it is final, since membership in the cone is NP-complete for $d\ge3$ (\cref{cor:hard-np}) and the cone therefore has no extended formulation of polynomial size unless $\mathrm P=\mathrm{NP}$ (\cref{cor:extended}). For $d=2$ such a formulation exists, the flow linear program of \cref{thm:lp}.
\item \emph{Existence, not uniqueness.} With finite data the distribution $\mu$ is never unique, and neither are the masses of the cells. What replaces uniqueness is stability: the criterion holds on an open set of prices (\cref{lem:open}), the finite counterpart of the fact that the profit function on an open cone of prices determines $\mu$ \citep[Thm.~2.1]{AgaltsovMolchanovShananin2018}.
\item \emph{Prices as the free variable.} In the identification problem of the school the data are fixed and the parameter of $h$ is chosen so that the moment problem becomes solvable. Here $h$ is free from the start, so the only quantities left to adjust are the prices themselves, and the least price correction that restores solvability (\cref{sec:lp}) is the natural measure of the distance of the data from the model.
\end{itemize}

\paragraph{Results.}
The first result (\cref{thm:solvability}) is the criterion. Call $S\subset[T]$ a \emph{reachable spectrum} at $P=(p_t)_{t\in[T]}$ if every $t\in S$ has a nonzero $\xi\in\Rp[d]$ with $\ip\xi{p_s-p_t}>0$ for all $s\notin S$. Then $y$ is generated by \eqref{eq:intro-output} for some neoclassical $h$ if and only if $y\in\cone\set{\ind S\mid S\text{ reachable}}$, and the same $h$ works for all prices near $P$. Prices on the boundary of the orthant and prices that coincide require a closure of the notion, \emph{weak solvability}, with the same criterion for weak witnesses (\cref{thm:closure}); its proof needs an uncrossing argument for coinciding prices (\cref{lem:monotone}).

The second result is a dichotomy in the number of inputs. For $d=2$ the reachable spectra are the sets cut out by convex chains of price points, and solvability is a flow linear program in an acyclic graph with $O(T^2)$ nodes (\cref{thm:lp}, \cref{cor:poly}); a certificate of unsolvability is a weight vector checked by one shortest-path computation (\cref{cor:farkas}). For $d\ge3$ both solvability and weak solvability are NP-complete (\cref{cor:hard-np}). The reduction from 3-colorability places prices on a paraboloid and turns solvability into membership in the stable set polytope of a graph drawn on a lattice; the drawing may have crossing edges, which is what makes the gadget work in three dimensions (\cref{sec:hardness}). The contrast with the fixed-$h$ regime is instructive. For a fixed two-input CES cost function the cells are polynomially many and the test is explicit; for an unknown $h$ the reachable spectra are exponentially many, and polynomiality at $d=2$ rests on their implicit description by chains. Stretchability plays no role here, because no particular arrangement has to be realized: the family of reachable spectra is quantified over all cost functions at once.

The third group of results concerns data that are not solvable. Prices are the least reliable part of industry statistics, and the natural question is the least relative correction of the prices that makes the outputs weakly solvable, measured by a weighted $\ell_r$ norm of the relative deviations. The box theorem (\cref{thm:corners}) reduces the question whether some prices in given boxes make $y$ weakly solvable to one query of the same oracle on the $2T$ box corners, with inequalities in place of equalities. The oracle is exact in every dimension: weakly reachable spectra are exactly the $0/1$ vectors satisfying \emph{cover inequalities} $x_t\le\sum_{s\in A}x_s$ with $|A|\le d+1$ (\cref{lem:cover}), and column generation over them, with all geometry checked in rational arithmetic, returns either a solution or a Farkas certificate together with the covers used (\cref{prop:oracle}). The least perturbation is then a monotonic optimization problem over box half-widths (\cref{thm:reduction}), and the certificates of the oracle yield combinatorial Benders cuts that are valid for every box in which the same covers are present (\cref{thm:certificate}). The resulting method keeps certified lower and upper bounds and terminates (\cref{thm:algo}); the cover cuts are what makes it converge on real data (\cref{sec:comp}).

\paragraph{Related work outside the school.}
In the econometric literature the Houthakker--Johansen model appears as the short-run production function built from microdata \citep{Hildenbrand1981}, where the distribution of capacities is observed. Nonparametric production analysis in the sense of \citet{Afriat1972}, \citet{HanochRothschild1972} and \citet{Varian1984} tests firm-level input and output quantities against a common technology; here quantities at the level of units are unobserved, and the common object is the unit cost function $h$ of a unit, not the production set of the industry. Measurement error in such tests was introduced by \citet{Varian1985} for quantities, and \citet{KuosmanenPostScholtes2007} perturb all observed quantities, choosing the $\ell_\infty$ norm for computability; Varian noted that errors in prices would make the constraints nonlinear and the problem nonstandard. The present problem is of that kind: the admissible prices form a nonconvex set whose membership test is NP-hard for $d\ge3$, and the contribution of \cref{sec:corners,sec:lp} is an exact method for it. Price perturbation for revealed preference tests of consumer demand, where the test is linear in prices, is a single mixed-integer program \citep{DemuynckRehbeck2023}; goodness-of-fit indices with one scalar per observation are surveyed in \citet{Cherchye2014}, and minimal-cost violation measures are treated by \citet{DeanMartin2016}. Criteria of the form ``the data lie in the cone of $0/1$ types'' are also the basis of the nonparametric analysis of random utility models \citep{McFaddenRichter1990,KitamuraStoye2018}, where the types are the cells of an arrangement of budget sets and the test is solved by column generation \citep{SmeuldersCherchyeDeRock2021}; the heterogeneity there is maximal, every agent has its own preference, whereas here all units share one $h$ and differ by their technology $x$. Complexity results for rationalizability of choice data go back to \citet{ApesteguiaBallester2010}; the dichotomy here is in the dimension of the inputs rather than in the structure of the data.

\paragraph{Formal verification and outline.}
The structural results are verified in Lean~4 with Mathlib; \cref{sec:lean} states exactly what is verified and what is taken from the literature, and the e-companion lists the Lean declarations for every numbered statement. \Cref{sec:solvability,sec:weak} prove the two criteria, \cref{sec:planar,sec:hardness} the dichotomy, \cref{sec:corners} the box theorem and the exact oracle, \cref{sec:lp} the perturbation algorithm, and \cref{sec:comp} reports computations on EU KLEMS and BEA--BLS data.
